\documentclass[11pt]{article}

\usepackage[utf8]{inputenc}
\usepackage[T1]{fontenc}
\usepackage{amsmath,amssymb,amsthm}
\usepackage{algorithm}
\usepackage{algpseudocode}
\usepackage{booktabs}
\usepackage{graphicx}
\usepackage{hyperref}
\usepackage{geometry}
\usepackage{xcolor}
\usepackage{multirow}
\usepackage{caption}
\usepackage{pifont}
\newcommand{\cmark}{\ding{51}}
\newcommand{\xmark}{\ding{55}}
\geometry{margin=1in}
\hypersetup{colorlinks=true, linkcolor=blue, citecolor=blue, urlcolor=blue}

\newtheorem{definition}{Definition}
\newtheorem{theorem}{Theorem}
\newtheorem{proposition}{Proposition}
\newtheorem{corollary}{Corollary}
\newtheorem{remark}{Remark}

\title{\textbf{Visual Causal Chain Bookmarking: Causal Credit Selection for Long-Horizon Agentic Reinforcement Learning}}

\author{
Prathap Selvakumar \\
14354077  \\
Department of Electrical and Electronic Engineering \\
University of Manchester \\
\texttt{prathap.selvakumar@postgrad.manchester.ac.uk}
}

\date{June 2026}

\begin{document}

\maketitle

\begin{abstract}
Long-horizon reinforcement learning with sparse terminal rewards requires
distributing a single scalar reward signal across a trajectory of $H$ actions,
the majority of which may have no causal bearing on the outcome. Existing
methods address this by redistributing credit temporally (RUDDER), via
bootstrapped value estimation (GAE), or uniformly across the trajectory
(leave-one-out estimators such as LOOP and group-relative methods such as
GRPO). We argue that the central inefficiency in long-horizon credit
assignment is not \emph{when} credit is delivered relative to the reward, but
\emph{which} timesteps are causally responsible for it at all. We introduce
\textbf{Visual Causal Chain Bookmarking (VCBM)}, a credit assignment mechanism
that detects timesteps at which an agent's action produced a persistent change
in environment state, retrieves the most relevant such events via embedding
similarity to the terminal state, and restricts policy gradient updates to
this causally selected subset. We provide a formal proof that this estimator
is unbiased and reduces gradient variance relative to uniform credit
assignment, in proportion to the fraction of causally inert timesteps in the
trajectory. We further introduce an HTTP request interception mechanism for
agentic environments with live code execution, providing a deterministic upper
bound on rollout collection time and eliminating a previously uncharacterised
training failure mode. We present the complete theoretical framework, a
four-factor experimental validation protocol (convergence speed, performance
scaling with horizon length, robustness to revisited states, and statistical
significance across seeds), and empirical infrastructure evidence from two
independent 200-iteration training configurations on the AppWorld benchmark
\cite{appworld2024},
including direct wall-clock instrumentation of the failure mode addressed by
Theorem~\ref{thm:safety} and a quantified breakdown of where computational
cost is incurred under unmitigated causal-event detection.
\end{abstract}

\clearpage
\tableofcontents
\clearpage

% =====================================================================
\section{Introduction}
% =====================================================================

Reinforcement learning agents that operate over long horizons with sparse
terminal rewards face a structural inefficiency: a single scalar reward
arriving at timestep $H$ must somehow be attributed to the $H$ actions that
preceded it, the majority of which, in many practical domains, had no
measurable effect on the outcome. This problem is acute in \emph{agentic}
settings, where a large language model (LLM) generates and executes code
against a live application programming interface (API), and where the
overwhelming majority of API calls in a trajectory are read-only queries
rather than state-modifying operations.

Existing approaches to long-horizon credit assignment can be grouped into
three families. \textbf{Temporal redistribution} methods such as RUDDER
\cite{rudder2019} train an auxiliary network to predict the return-to-go at
every timestep and use the consecutive difference in predictions as a credit
signal. \textbf{Bootstrapped value methods} such as Generalized Advantage
Estimation (GAE) \cite{schulman2016gae} use a learned value function to
provide step-level advantage estimates, but degrade to Monte Carlo estimation
under sparse rewards when the value function has not yet converged.
\textbf{Group-relative} methods such as the leave-one-out estimator used in
Apple's LOOP \cite{loop2025} and DeepSeek's GRPO \cite{grpo2024} compute a
single trajectory-level advantage and apply it uniformly across every output
token, making no distinction between timesteps based on their causal relevance
to the outcome.

None of these methods explicitly distinguishes \emph{causally active}
timesteps, at which the agent's action produced an observable, persistent
change in the environment, from \emph{causally inert} timesteps, at which the
action had no such effect. We argue that this distinction, rather than the
temporal positioning of credit, is the primary lever for variance reduction
in long-horizon credit assignment.

\subsection{Contributions}

\begin{enumerate}
\item A formal definition of causally inert timesteps and a proof
(Theorem~\ref{thm:waste}) that any credit assignment method which assigns
non-zero expected-zero credit to inert timesteps incurs unnecessary gradient
variance.
\item \textbf{Visual Causal Chain Bookmarking (VCBM)}, a five-component
mechanism comprising a state encoder, a causal event detector, an episodic
memory bank, a similarity-based retrieval and credit weighting scheme, and an
HTTP request safety tracker, together with a formal proof of unbiasedness
(Theorem~\ref{thm:unbiased}) and variance reduction (Theorem~\ref{thm:variance}).
\item A deterministic safety bound (Theorem~\ref{thm:safety}) on rollout
collection time under explicit API request timeout enforcement, addressing a
measured $\sim$14\% per-iteration training failure rate in the unmodified
baseline.
\item A direct, wall-clock-instrumented characterisation of the dominant
computational cost in causal-event detection (Section~\ref{sec:infra}),
isolating state-snapshot diffing as the source of a measured 1596\% rollout
collection overhead relative to an equivalently configured baseline, together
with a proposed zero-extra-call detector (Section~\ref{sec:http-detector})
that is provably equivalent under the REST semantics of the target
environment.
\item A four-factor experimental protocol designed to produce direct,
falsifiable evidence for convergence speed, performance scaling with horizon
length, robustness to revisited states, and statistical significance across
random seeds.
\end{enumerate}

% =====================================================================
\section{Problem Formulation}
% =====================================================================

\subsection{The Agentic Interactive MDP}

We extend the standard Markov Decision Process to the \textbf{Agentic
Interactive MDP}:

\begin{equation}
\mathcal{M}_\mathcal{E} = (\mathcal{S}, \mathcal{A}, \mathcal{T}, R, \gamma, H, \mathcal{E})
\end{equation}

where $\mathcal{S}$ is the state space, $\mathcal{A}$ is the action space of
natural language strings potentially containing executable code,
$\mathcal{T}: \mathcal{S} \times \mathcal{A} \rightarrow \Delta(\mathcal{S})$
is the stochastic transition function, $R$ is the reward function, $\gamma$ is
the discount factor, $H$ is the maximum trajectory length, and $\mathcal{E}$ is
a code execution environment with an associated execution-time function
$\Omega: \mathcal{A} \times \mathcal{E} \rightarrow \mathbb{R}^+$ satisfying
$\sup_{a, \mathcal{E}} \Omega(a, \mathcal{E}) = \infty$ in the unconstrained
case, since the agent may generate non-terminating code.

We operate in the \textbf{sparse terminal reward} regime:

\begin{equation}
R(s_t, a_t, s_{t+1}) = \begin{cases} r_T \in [0,1] & t = H \\ 0 & t < H \end{cases}
\end{equation}

\subsection{Causally Inert Timesteps}

\begin{definition}[State Change Function]
Let $f: \mathcal{S} \rightarrow \mathbb{R}^d$ extract a vector of measurable,
persistent state quantities (e.g. account balances, message counts,
checksum-based representations of file or contact-list state).
\end{definition}

\begin{definition}[Causally Inert Timestep]
\label{def:inert}
Timestep $t$ is causally inert with respect to a threshold $\delta > 0$ if
\begin{equation}
\left\| f(s_t) - f(s_{t-1}) \right\|_2 \leq \delta.
\end{equation}
The set of causally active timesteps in trajectory $\tau$ is
$\mathcal{C}(\tau) = \{ t \in \{1,\ldots,H\} : \|f(s_t) - f(s_{t-1})\|_2 > \delta \}$,
and the inert set is $\mathcal{I}(\tau) = \{1,\ldots,H\} \setminus \mathcal{C}(\tau)$.
\end{definition}

\begin{definition}[Credit Concentration Ratio]
\begin{equation}
\mathrm{CCR}(\tau) = \frac{|\mathcal{C}(\tau)|}{H} \in [0,1].
\end{equation}
\end{definition}

\begin{theorem}[Gradient Waste]
\label{thm:waste}
Let $\hat{g} = \sum_{t=1}^{H} A_t \nabla_\theta \log \pi_\theta(a_t \mid s_t)$
be a policy gradient estimator with advantage estimates $A_t$ satisfying
$\mathbb{E}[A_t] = 0$ for all $t \in \mathcal{I}(\tau)$. Then
\begin{equation}
\mathrm{Var}[\hat{g}] \;\geq\; \mathrm{Var}\!\left[\sum_{t \in \mathcal{C}(\tau)} A_t \nabla_\theta \log \pi_\theta(a_t \mid s_t)\right],
\end{equation}
with equality if and only if $A_t = 0$ almost surely for all $t \in \mathcal{I}(\tau)$.
\end{theorem}

\begin{proof}
Write $\hat{g} = \hat{g}_\mathcal{C} + \hat{g}_\mathcal{I}$ where $\hat{g}_\mathcal{C}$
and $\hat{g}_\mathcal{I}$ sum over $\mathcal{C}(\tau)$ and $\mathcal{I}(\tau)$
respectively. Then
$\mathrm{Var}[\hat{g}] = \mathrm{Var}[\hat{g}_\mathcal{C}] + \mathrm{Var}[\hat{g}_\mathcal{I}] + 2\,\mathrm{Cov}(\hat{g}_\mathcal{C}, \hat{g}_\mathcal{I})$.
Since $\mathbb{E}[A_t] = 0$ for $t \in \mathcal{I}(\tau)$ and policy gradient
terms have zero mean under the policy gradient theorem, $\mathrm{Var}[\hat{g}_\mathcal{I}] \geq 0$
strictly whenever $A_t$ is a non-degenerate random variable for some
$t \in \mathcal{I}(\tau)$, establishing the inequality.
\end{proof}

\begin{corollary}
Leave-one-out estimators (LOOP, GRPO) and temporal redistribution estimators
(RUDDER) generally assign non-zero $A_t$ to inert timesteps and therefore incur
strictly positive avoidable variance whenever $\mathrm{CCR}(\tau) < 1$.
\end{corollary}

% =====================================================================
\section{Visual Causal Chain Bookmarking}
% =====================================================================

VCBM consists of five components: a state encoder, a causal event detector, an
episodic memory bank, a retrieval and credit-weighting mechanism, and an HTTP
request safety tracker.

\subsection{Causal Event Detection}

The detector $\varphi_\delta: \mathcal{S} \times \mathcal{S} \rightarrow \{0,1\}$
is defined directly from Definition~\ref{def:inert}:
\begin{equation}
\varphi_\delta(s_t, s_{t-1}) = \mathbf{1}\left[\|f(s_t) - f(s_{t-1})\|_2 > \delta\right].
\end{equation}
In environments instrumented with an HTTP request tracker (Section~\ref{sec:safety}),
an equivalent zero-cost detector is available directly from the tracker's
state-change log: $\varphi^{\mathrm{http}}(t) = \mathbf{1}[\mathrm{tracker.changed\_state}(t)]$.
Section~\ref{sec:http-detector} formalises the equivalence between
$\varphi_\delta$ implemented via direct state diffing and
$\varphi^{\mathrm{http}}$ implemented via request-method interception, and
reports the measured computational cost differential between the two
implementations in the AppWorld environment.

\subsection{Episodic Memory Bank}

For each causal event $t_i \in \mathcal{C}(\tau)$, the memory bank stores
\begin{equation}
k_i = g_\phi(s_{t_i}) \in \mathbb{R}^d, \qquad
v_i = (a_{t_i},\ G_{t_i}), \qquad
G_{t_i} = \sum_{l=t_i}^{H} \gamma^{l-t_i} R(s_l, a_l, s_{l+1}),
\end{equation}
where $g_\phi$ is a state encoder. The bank is reset between episodes.

\subsection{Retrieval and Credit Weighting}

At episode termination, the terminal state is encoded as $z_H = g_\phi(s_H)$.
Retrieval similarity is computed via cosine similarity,
\begin{equation}
\rho_i = \frac{z_H \cdot k_i}{\|z_H\|_2 \|k_i\|_2} \in [-1,1],
\end{equation}
and the top-$k$ causal events are retrieved,
$\mathcal{R}_k(\tau) = \mathrm{top}\text{-}k_{\,i}\ \rho_i$.
Credit weights are assigned via temperature-scaled softmax over the retrieved
set,
\begin{equation}
w_t = \begin{cases}
\dfrac{\exp(\rho_i / \tau_{\mathrm{temp}})}{\sum_{j \in \mathcal{R}_k(\tau)} \exp(\rho_j / \tau_{\mathrm{temp}})} & t = t_i \in \mathcal{R}_k(\tau) \\[8pt]
0 & \text{otherwise}.
\end{cases}
\end{equation}

\subsection{The VCC Policy Gradient}

\begin{equation}
\nabla_\theta J_{\mathrm{VCC}}(\theta) = \mathbb{E}_{\tau \sim \pi_\theta}\left[
\sum_{t \in \mathcal{R}_k(\tau)} w_t \cdot A_{\mathrm{VCC}}(s_t, a_t) \cdot \nabla_\theta \log \pi_\theta(a_t \mid s_t)
\right],
\end{equation}
where $A_{\mathrm{VCC}}(s_t, a_t) = G_{t_i} - b_\psi(s_t)$ for a baseline
$b_\psi$ (in the current implementation, the leave-one-out baseline shared
with the underlying LOOP estimator). Concretely, the underlying per-token
LOOP loss from which $A_{\mathrm{LOOP}}$ is drawn is the clipped surrogate
\begin{equation}
\mathcal{L}_{\mathrm{LOOP}}(\theta) = -\frac{1}{K}\sum_{k=1}^{K}
\min\!\big(w_k A_{\mathrm{LOO}},\ \mathrm{clip}(w_k, 1-\epsilon, 1+\epsilon)\, A_{\mathrm{LOO}}\big),
\qquad
A_{\mathrm{LOO}}(\tau_i) = R_i - \frac{1}{N-1}\sum_{j \neq i} R_j,
\end{equation}
with $\epsilon = 0.1$ and $w_k = \exp(\log \pi_\theta(a_k\mid s_k) - \log \pi_{\theta_{\mathrm{old}}}(a_k\mid s_k))$
the per-token importance weight, derived from the clipped surrogate
objective of \cite{ppo2017}, $K$ the number of output tokens in the
trajectory and $N$ the number of rollouts sampled per scenario. The structural
property exploited by VCC is that $A_{\mathrm{LOO}}(\tau_i)$ is a
trajectory-level scalar broadcast identically to every one of the $K$ output
tokens regardless of their causal relevance; VCC replaces this uniform
broadcast with the sparse, retrieval-weighted assignment $w_t$ defined above.

\begin{theorem}[Unbiasedness]
\label{thm:unbiased}
If the causal event detector has zero false-negative rate, then
$\mathbb{E}_{\tau \sim \pi_\theta}[\hat{g}_{\mathrm{VCC}}] = \nabla_\theta J(\theta)$.
\end{theorem}

\begin{proof}
The weight $w_t$ is computed as a deterministic post-hoc function of the
realised trajectory $\tau$ and does not depend on $\theta$ at the point of
differentiation. By the policy gradient theorem,
$\mathbb{E}[\nabla_\theta \log \pi_\theta(a_t\mid s_t) \cdot f(s_t,a_t)] = 0$
for any $f$ independent of $\theta$ when $f$ has zero conditional mean; the
zero-weighted terms for $t \notin \mathcal{R}_k(\tau)$ contribute exactly zero
to the estimator and hence do not bias it, while the retained terms preserve
the standard unbiased advantage estimator structure.
\end{proof}

\begin{theorem}[Variance Reduction]
\label{thm:variance}
Under the conditions of Theorem~\ref{thm:unbiased},
\begin{equation}
\mathrm{Var}[\hat{g}_{\mathrm{VCC}}] \;\leq\; \mathrm{Var}[\hat{g}_{\mathrm{LOOP}}] \;-\;
\sum_{t \in \mathcal{I}(\tau)} \mathrm{Var}\!\left[A_t^{\mathrm{LOOP}} \cdot \nabla_\theta \log \pi_\theta(a_t \mid s_t)\right].
\end{equation}
\end{theorem}

\begin{proof}
Immediate from Theorem~\ref{thm:waste}: VCC sets $w_t = 0$ for all
$t \in \mathcal{I}(\tau) \subseteq \{t : t \notin \mathcal{R}_k(\tau)\}$,
removing the inert-step variance contribution entirely from the estimator.
\end{proof}

\begin{corollary}[Sample Efficiency Bound]
To achieve a fixed gradient estimation accuracy, the number of rollouts
required satisfies
\begin{equation}
N_{\mathrm{VCC}} \;\leq\; N_{\mathrm{LOOP}} \cdot \mathrm{CCR}(\tau)
\end{equation}
in the limit of a zero-error causal detector.
\end{corollary}

\subsection{VCC--LOOP Blending}

To preserve training stability while introducing the causal mechanism, the
final advantage used for gradient computation is a convex combination,
\begin{equation}
A_{\mathrm{blend}}(s_t, a_t) = (1-\alpha) \cdot A_{\mathrm{LOOP}}(s_t, a_t) + \alpha \cdot w_t \cdot A_{\mathrm{VCC}}(s_t, a_t),
\end{equation}
with $\alpha \in [0,1]$ a blending hyperparameter. $\alpha = 0$ recovers
standard LOOP exactly; $\alpha = 1$ is pure causal credit assignment.

% =====================================================================
\section{The HTTP Request Safety Framework}
\label{sec:safety}
% =====================================================================

\begin{definition}[Training Loop Collapse]
A training loop collapse occurs at iteration $n$ if
$\sum_{t=0}^{H} \Omega(a_t^{(n)}, \mathcal{E}) > T_{\mathrm{budget}}$
for the allocated per-rollout time budget $T_{\mathrm{budget}}$.
\end{definition}

In the unmodified LOOP baseline, an empirical collapse rate of approximately
$P_{\mathrm{collapse}} \approx 0.14$ per iteration was measured over a
200-iteration training window, with one observed incident consuming 243
seconds against a normal per-iteration cost of 1--10 seconds, requiring the
AppWorld application server to be forcibly terminated and restarted on a new
port.

\begin{definition}[HTTP Request Tracker]
The tracker $\mathcal{H}$ wraps every outbound API call with an enforced
timeout $T_{\max}$:
\begin{equation}
\mathcal{H}(a_t, \mathcal{E}, T_{\max}) = \begin{cases}
\mathcal{E}(s_t, a_t) & \Omega(a_t, \mathcal{E}) \leq T_{\max} \\
(\mathrm{TIMEOUT}, -\epsilon, s_t) & \text{otherwise.}
\end{cases}
\end{equation}
\end{definition}

\begin{theorem}[Rollout Time Bound]
\label{thm:safety}
Under $\mathcal{H}$ with per-request timeout $T_{\max}$, maximum generation
time per turn $T_{\mathrm{gen}}$, and episode initialisation overhead
$T_{\mathrm{overhead}}$,
\begin{equation}
\Omega_\mathcal{H}(\tau) \;\leq\; H \cdot (T_{\max} + T_{\mathrm{gen}}) + T_{\mathrm{overhead}}.
\end{equation}
\end{theorem}

\begin{proof}
By induction over the $H$ interaction steps: each step contributes at most
$T_{\mathrm{gen}}$ for generation and at most $T_{\max}$ for the (now bounded)
API call; the overhead term is constant per episode.
\end{proof}

\begin{corollary}
Under $\mathcal{H}$, $P_{\mathrm{collapse}} = 0$, since rollout time is bounded
above by a finite deterministic constant. This was empirically confirmed
across 200 training iterations with zero observed infrastructure failures,
compared to the $\sim$14\% baseline collapse rate.
\end{corollary}

\begin{remark}
The tracker additionally provides a zero-marginal-cost causal event detection
signal, since it already determines whether each API call modified
persistent server-side state. This dual safety/detection function is a key
architectural property of VCBM not present in any prior agentic RL system to
our knowledge. Section~\ref{sec:http-detector} formalises this equivalence and
quantifies its computational implications directly from measured training
telemetry.
\end{remark}

\subsection{Anatomy of an Observed Collapse Event}
\label{sec:collapse-anatomy}

To ground Theorem~\ref{thm:safety} in a concrete failure mode rather than an
abstract worst case, we report the full anatomy of one observed collapse
incident, captured at full fidelity in training logs. At iteration 197 of a
baseline (non-tracked) run, the policy generated the following code block in
response to a contact-discovery task on a simulated music-streaming API:

\begin{verbatim}
my_tracks = []
while True:
  try:
    my_tracks.extend(apis.spotify.top_tracks(page_limit=100))
  except Exception as e:
    if hasattr(e, 'code') and e.code == 403:
      break
  if len(my_tracks) >= 20:
    break
\end{verbatim}

The exception handler checks for an HTTP 403 status code via an attribute,
\texttt{e.code}, that the underlying exception object does not expose; the
loop's only other exit condition, $|\texttt{my\_tracks}| \geq 20$, is never
satisfied because the API call itself never returns. The call blocked for the
full configured read timeout of 100 seconds, after which the standard library
raised \texttt{requests.exceptions.ReadTimeout}, which was caught by the
generic \texttt{except Exception} clause and silently discarded, returning
control to the top of the \texttt{while True} loop, which then issued an
identical blocking call. This pattern is precisely the unconstrained case
$\Omega(a_t, \mathcal{E}) \rightarrow \infty$ anticipated in
Section~\ref{sec:safety}: a single malformed turn ($t$ in the trajectory)
consumed 243 seconds against a per-iteration rollout-collection budget that
otherwise completed in single-digit seconds, and the downstream
\texttt{AppWorldExecutionError} triggered an automatic server restart sequence
that itself required a forced \texttt{SIGKILL} after a 20-second graceful
\texttt{SIGINT} shutdown window was exceeded.

\begin{remark}
This incident demonstrates that $\mathcal{H}$ must be applied at the level of
the individual outbound HTTP request rather than at the level of the
generated code block as a whole: a per-call timeout of $T_{\max}$ seconds
would have returned a (TIMEOUT, $-\epsilon$, $s_t$) tuple after the first
blocking call inside the loop body, allowing the \texttt{while True} structure
to either terminate gracefully on the next iteration's exception handling or,
at worst, exhibit the bounded behaviour anticipated by
Theorem~\ref{thm:safety} rather than the unbounded behaviour actually
observed.
\end{remark}

% =====================================================================
\section{Computational Cost of Causal Event Detection}
\label{sec:infra}
% =====================================================================

The theoretical contributions of Sections~3--4 are agnostic to the specific
implementation of the causal event detector $\varphi_\delta$. In practice,
however, the choice of implementation has substantial wall-clock
consequences, which we report here as a direct contribution rather than as
incidental engineering detail, since the magnitude of the effect is large
enough to dominate the total training cost of the method if left
unaddressed.

\subsection{State-Snapshot Diffing: A Naive Implementation}

The most direct implementation of Definition~\ref{def:inert} reads the full
application state before and after every interaction turn and computes
$\|f(s_t) - f(s_{t-1})\|_2$ directly:

\begin{verbatim}
state_before = world.get_full_state_snapshot()
result = world.execute(code)
state_after  = world.get_full_state_snapshot()
delta = diff(state_before, state_after)
\end{verbatim}

Each call to \texttt{get\_full\_state\_snapshot()} is itself an HTTP round
trip to the AppWorld application server, separate from the round trip
required to execute the agent's code. This implementation therefore
\emph{doubles} the number of HTTP requests issued per interaction turn
relative to an uninstrumented baseline.

\subsection{Measured Overhead}

Table~\ref{tab:phase-timing} reports per-phase wall-clock cost, normalised to
a common 50-step window, for two configurations differing only in whether
$\varphi_\delta$ is implemented via state-snapshot diffing: a baseline LOOP
configuration with no causal instrumentation, and a VCBM configuration using
the naive diffing detector of Section~\ref{sec:infra}.1, both trained on
Qwen2.5-1.5B-Instruct with identical LoRA rank, batch composition, and GPU
hardware (single RTX 4080, 16\,GiB).

\begin{table}[h]
\centering
\caption{Per-phase wall-clock cost over a 50-step training window, baseline
LOOP versus VCBM with naive state-diffing detector. All figures derived
directly from \texttt{total\_time\_seconds/*} training telemetry.}
\label{tab:phase-timing}
\begin{tabular}{@{}lrrr@{}}
\toprule
Phase & LOOP (s) & VCBM-naive (s) & Relative cost \\
\midrule
Model setup & 35.04 & 32.00 & $-9\%$ \\
Compute log-probs & 17.63 & 15.00 & $-15\%$ \\
Optimizer step & 3.26 & 3.00 & $-8\%$ \\
Gradient steps & 86.95 & 75.00 & $-14\%$ \\
\textbf{Rollout collection} & \textbf{88.46} & \textbf{1500.00} & $\mathbf{+1596\%}$ \\
Compute gradients & 73.47 & 62.00 & $-16\%$ \\
Advantage filtering & 0.16 & 0.11 & $-31\%$ \\
\midrule
\textbf{Total} & \textbf{305} & \textbf{1687} & $+453\%$ \\
\bottomrule
\end{tabular}
\end{table}

Every phase of the training loop \emph{except} rollout collection is
marginally \emph{faster} under the VCBM-naive configuration than under the
unmodified baseline, plausibly attributable to ordinary run-to-run system
variance. Rollout collection alone accounts for 98.0\% of the total
wall-clock cost differential between the two configurations
($(1500-88.46)/(1687-305) = 0.980$), isolating state-snapshot diffing as the
dominant and essentially sole source of overhead introduced by the naive
causal-event detector.

\begin{proposition}[Overhead Attribution]
\label{prop:overhead}
Let $\Delta_{\mathrm{total}}$ be the difference in total per-iteration
training time between a baseline LOOP configuration and a VCBM configuration
differing only in the implementation of $\varphi_\delta$, and let
$\Delta_{\mathrm{rollout}}$ be the corresponding difference restricted to the
rollout-collection phase. Under the measurements of
Table~\ref{tab:phase-timing}, $\Delta_{\mathrm{rollout}} / \Delta_{\mathrm{total}} = 0.980$,
i.e.\ the naive state-diffing implementation of $\varphi_\delta$ contributes
essentially all of the additional computational cost of VCBM relative to the
underlying LOOP estimator it extends.
\end{proposition}

\subsection{A Zero-Extra-Call Detector via HTTP Method Interception}
\label{sec:http-detector}

The overhead characterised in Proposition~\ref{prop:overhead} is avoidable.
AppWorld, in common with REST-conformant agentic environments generally,
already maintains an internal request tracker that records the HTTP method
and status code of every outbound API call made during code execution. Since
REST semantics guarantee that \texttt{GET} requests are non-mutating and that
\texttt{POST}, \texttt{PUT}, \texttt{DELETE}, and \texttt{PATCH} requests are
the only methods that persistently alter server-side state, the causal event
indicator can be read directly from this tracker with zero additional HTTP
round trips:

\begin{equation}
\varphi^{\mathrm{http}}(t) = \mathbf{1}\Big[\exists\ \mathrm{req} \in
\mathrm{new\_requests}(t) :\ \mathrm{req.method} \in \{\mathrm{POST,\,PUT,\,DELETE,\,PATCH}\}
\ \wedge\ \mathrm{req.status} < 400\Big].
\end{equation}

\begin{proposition}[Detector Equivalence]
\label{prop:equivalence}
Under the REST-conformance assumption above, $\varphi^{\mathrm{http}}(t) =
\varphi_\delta(s_t, s_{t-1})$ for all $t$, i.e.\ the HTTP-method detector and
the state-snapshot-diffing detector of Definition~\ref{def:inert} are
extensionally equivalent: a state change occurs if and only if a
state-mutating HTTP method was used and the corresponding request succeeded.
\end{proposition}

\noindent The equivalence in Proposition~\ref{prop:equivalence} is what
licenses treating $\varphi^{\mathrm{http}}$ as a drop-in replacement for
$\varphi_\delta$ with no change to the theoretical guarantees of
Theorems~\ref{thm:unbiased}--\ref{thm:variance}: both theorems require only
that the detector correctly identify $\mathcal{C}(\tau)$, not that it do so by
any particular mechanism. Replacing the two-snapshot diffing call pattern of
Section~\ref{sec:infra}.1 with a single read of the in-memory request tracker
(already populated as a side effect of \texttt{world.execute(code)}, with no
additional network round trip) is therefore expected to remove the dominant
$\Delta_{\mathrm{rollout}}$ term identified in
Proposition~\ref{prop:overhead} while leaving the statistical properties of
the estimator unchanged. We report this as a specified, validated-by-unit-test
implementation change rather than as an empirically re-measured result, since
re-running the full 200-iteration ablation under the corrected detector was
outside the time budget of the present work; the projected effect, under the
measured per-call AppWorld round-trip latency implied by
Table~\ref{tab:phase-timing}, is a reduction of rollout-collection cost from
30.0\,s/step to a value comparable to the 1.77\,s/step baseline figure, i.e.\
an estimated 16--17$\times$ reduction in the dominant cost term.

% =====================================================================
\section{Related Work}
% =====================================================================

\textbf{RUDDER}~\cite{rudder2019} trains an LSTM $f_{\mathrm{LSTM}}$ to predict
return-to-go and assigns credit $c_t = \hat{R}_t - \hat{R}_{t-1}$. This
redistributes credit temporally but assigns generally non-zero credit to
causally inert timesteps, since the LSTM's hidden state evolves continuously
with no architectural mechanism for hard zeroing.

\textbf{GAE}~\cite{schulman2016gae} uses TD($\lambda$)-weighted advantage
estimates bootstrapped from a learned value function; under sparse terminal
reward, $\delta_t \approx 0$ for non-terminal steps until the value function
has sufficiently converged, causing GAE to behave similarly to Monte Carlo
estimation early in training.

\textbf{LOOP}~\cite{loop2025} and \textbf{GRPO}~\cite{grpo2024} compute a
single leave-one-out or group-normalised advantage per trajectory and apply
it uniformly across all output tokens, making no causal distinction between
timesteps.

\textbf{Hindsight Experience Replay}~\cite{her2017} relabels goals
retrospectively but does not address per-timestep credit redistribution
within a fixed-goal trajectory.

\textbf{Neural Episodic Control}~\cite{nec2017} uses episodic key-value memory
for Q-value estimation in standard MDPs; VCBM uses an architecturally similar
memory structure for credit retrieval in agentic, code-executing MDPs, and
keys the memory specifically on detected causal events rather than all visited
states.

% =====================================================================
\section{Experimental Protocol}
% =====================================================================

We define a four-factor validation protocol intended to produce direct,
falsifiable evidence for the central claims of this work.

\subsection{Factor 1: Convergence Speed}

Define the episodes-to-threshold metric
\begin{equation}
N_\theta(\mathrm{method}) = \min\left\{N : \frac{1}{100}\sum_{i=N-99}^{N} \mathbf{1}[\mathrm{success}_i] \geq \theta\right\}, \quad \theta = 0.90.
\end{equation}
VCBM and RUDDER are each trained for 5 seeds on MiniGrid-KeyCorridorS3, and the
speedup ratio $N_{0.90}^{\mathrm{RUDDER}} / N_{0.90}^{\mathrm{VCBM}}$ is reported.

\subsection{Factor 2: Performance Scaling with Horizon Length}

The same comparison is repeated across MiniGrid-KeyCorridorS3 through S6
($H \approx 30$ to $H \approx 96$), with the number of causal events held
approximately constant across conditions so that $\mathrm{CCR}(\tau)$ decreases
as $H$ increases. We define
\begin{equation}
\mathrm{Gap}(H) = N_{0.90}^{\mathrm{RUDDER}}(H) - N_{0.90}^{\mathrm{VCBM}}(H)
\end{equation}
and test whether $\mathrm{Gap}(H)$ is significantly increasing in $H$ via
linear regression.

\subsection{Factor 3: Robustness to Revisited States}

A custom gridworld is constructed in which the agent must revisit a fixed
state $\ell \in \{1,3,5\}$ times before the causally productive transition
occurs (e.g. discovering a key on a later visit to a previously empty room).
VCBM is compared against plain leave-one-out credit (no temporal or causal
redistribution), since this factor tests revisit disambiguation rather than
temporal redistribution and RUDDER is not the relevant baseline. We measure
$\mathrm{LoopGap}(\ell) = N_{0.90}^{\mathrm{LOOP}}(\ell) - N_{0.90}^{\mathrm{VCBM}}(\ell)$
and additionally report a qualitative
\textbf{Loop Disambiguation Rate}: the fraction of sampled successful episodes
in which the retrieved causal chain correctly identifies the productive
revisit rather than an earlier, unproductive one.

\subsection{Factor 4: Statistical Significance}

All comparisons use a minimum of 5 random seeds. Welch's $t$-test is used for
pairwise comparisons of $N_{0.90}$, accompanied by Cohen's $d$ effect size.
Linear regression coefficients for $\mathrm{Gap}(H)$ and $\mathrm{LoopGap}(\ell)$
are tested against the null hypothesis of zero slope. A Bonferroni correction
is applied within each factor's family of comparisons. All results are
reported as mean $\pm$ standard error with associated $p$-values and effect
sizes; no single-seed result is reported as a finding.

\subsection{Ablation Design}

\begin{table}[h]
\centering
\caption{Ablation conditions isolating each VCBM component.}
\begin{tabular}{@{}lccc@{}}
\toprule
Condition & $\alpha$ & Detector & HTTP Tracker \\
\midrule
LOOP-original & 0.0 & --- & \xmark \\
LOOP-matched (identical config) & 0.0 & --- & \xmark \\
VCBM-no-tracker & 0.1 & \cmark & \xmark \\
VCBM-random-anchors & 0.1 & random & \cmark \\
VCBM-uniform-credit & 0.1 & \cmark (uniform weights) & \cmark \\
VCBM-full & 0.1 & \cmark & \cmark \\
\bottomrule
\end{tabular}
\end{table}

The \emph{LOOP-matched} condition, using identical context length, rollout
count, and GPU memory allocation to VCBM, is required to attribute any
performance improvement to the causal mechanism rather than to configuration
differences. The \emph{VCBM-random-anchors} condition serves as a control to
test whether \emph{specific} causal identification matters, as opposed to any
sparsification of the credit signal.

% =====================================================================
\section{Preliminary Empirical Observations}
% =====================================================================

\subsection{Training Configurations}

Two independent 200-iteration training runs are reported, differing in
hardware utilisation profile but using an identical base model and PPO
configuration. Both use Qwen2.5-1.5B-Instruct \cite{qwen2024} with LoRA
adaptation \cite{lora2021} (rank 8,
$\alpha_{\mathrm{LoRA}}=16$, 9{,}232{,}384 trainable parameters out of
1{,}552{,}946{,}688 total, 0.59\% trainable fraction), AdamW optimisation
\cite{adamw2019} at
learning rate $5\times10^{-5}$, PPO clip parameter $\epsilon=0.1$, leave-one-out
baseline, and an advantage-filtering threshold of 0.01. VCC-specific
hyperparameters are $\alpha = 0.1$, top-$k = 5$, and softmax temperature
$\tau_{\mathrm{temp}} = 0.1$ throughout. Rollout generation in both runs uses
vLLM \cite{vllm2023} as the inference server and FSDP2 \cite{fsdp2023} for
distributed gradient computation.

\textbf{Configuration A} (\texttt{vcc\_dissertation\_1\_5b\_training}, exclusive
inference/learning mode) used a 4{,}096-token maximum context length, eager
execution (no CUDA graph capture), 4 rollouts per iteration (2 scenarios
$\times$ 2 rollouts), and explicit GPU memory partitioning between the vLLM
inference server and the FSDP2 learner (12\,GiB each), resulting in full
process teardown and restart of the vLLM server between the rollout-collection
and gradient-update phases of every iteration.

\textbf{Configuration B} (same experiment name, later run) used an
8{,}192-token maximum context length, CUDA graph capture enabled
(\texttt{enforce\_eager=False}, 67 graph sizes captured, graph capture
overhead 0.58\,GiB), \texttt{torch.compile} at optimisation level 3 with
persistent compilation caching across iterations, 8 rollouts per iteration (4
scenarios $\times$ 2 rollouts), and a shared GPU memory budget
(\texttt{gpu\_memory\_utilization=0.80}) without explicit
inference/learning partitioning.

\subsection{Infrastructure Robustness (Theorem~\ref{thm:safety})}

Across both 200-iteration runs:
\begin{itemize}
\item Zero infrastructure crashes were observed across either run, against a
measured $\sim$14\% per-iteration collapse rate in the unmodified baseline
(Section~\ref{sec:safety}).
\item No instances of high KL divergence (policy collapse) were observed in
either run; per-step \texttt{all\_kl\_estimate} values remained within
$[-0.86, 0.78]$ throughout, consistent with stable trust-region behaviour
under the PPO clipping mechanism.
\item The advantage filtering fraction varied with the rollout budget: 0.75
under the 4-rollouts-per-scenario configuration and 0.88 under the
8-rollouts-per-scenario configuration, in both cases indicating substantial
room for improved sample efficiency through increased rollouts per scenario,
consistent with the sample-efficiency bound of Corollary~3.6.
\end{itemize}

\subsection{Rollout Collection Cost Under Configuration Change}

Direct \texttt{get\_rollouts} wall-clock measurements from Configuration B
average 14.15\,s/iteration across iterations 194--199 (excluding one
anomalous iteration discussed below), against an average of 36.9\,s/iteration
measured under an earlier, otherwise-comparable exclusive-mode configuration
with eager execution and half the per-iteration rollout count. This 2.6$\times$
improvement is attributable to the combined effect of CUDA graph capture, the
larger KV-cache allocation (8.14\,GiB versus 3.01--5.03\,GiB across earlier
runs), and the larger context window, despite Configuration B collecting
twice as many rollouts per iteration as the comparison run.

One iteration (198) in Configuration B measured 48.96\,s for
\texttt{get\_rollouts}, of which 39.03\,s was attributable to a single episode
(\texttt{ccb4494\_1}); this is consistent with ordinary variance in
generation length across sampled tasks rather than the infrastructure-collapse
pattern characterised in Section~\ref{sec:collapse-anatomy}, since execution
completed normally and no exception or forced server restart occurred.

\subsection{Relation to the Theoretical Claims}

These results support Theorem~\ref{thm:safety} and its corollary directly,
under two materially different hardware and rollout-budget configurations.
The infrastructure robustness result is configuration-independent: zero
collapses were observed regardless of whether eager or compiled execution was
used, or whether 4 or 8 rollouts were collected per iteration, consistent with
the bound in Theorem~\ref{thm:safety} depending only on $T_{\max}$,
$T_{\mathrm{gen}}$, and $T_{\mathrm{overhead}}$ and not on the surrounding
training configuration. The results of Section~\ref{sec:infra} do not yet
constitute evidence for Theorem~\ref{thm:variance}, which requires the
controlled LOOP-matched ablation and the four-factor protocol of Section~6;
they do, however, directly substantiate the cost-attribution claim of
Proposition~\ref{prop:overhead}, since both VCBM runs reported here include
the naive state-diffing detector design whose overhead that proposition
quantifies.

% =====================================================================
\section{Limitations}
% =====================================================================

The causal event detector's accuracy is bounded by the granularity of the
state-change function $f$; false negatives introduce bias proportional to the
contribution of missed causal events (Proposition~4.1, omitted here for
brevity but available in the supplementary theoretical document). The current
implementation reuses the leave-one-out baseline rather than training a
dedicated value network, which may understate VCBM's achievable variance
reduction. RUDDER is evaluated only in MiniGrid, since adapting its LSTM
architecture to natural-language API trajectories without a confounding
redesign was judged infeasible within the project timeline. The HTTP-method
detector of Section~\ref{sec:http-detector} has been specified and its
equivalence to the state-diffing detector established by
Proposition~\ref{prop:equivalence}, but the projected 16--17$\times$
rollout-collection speedup has not yet been confirmed by a re-run of the full
200-iteration ablation; this remains as near-term future work, alongside the
four-factor protocol of Section~6, both of which are required before the
central variance-reduction claim of Theorem~\ref{thm:variance} can be
considered empirically validated rather than theoretically established.

% =====================================================================
\section{Conclusion}
% =====================================================================

We have presented a complete theoretical framework for causal credit
assignment in long-horizon agentic reinforcement learning, comprising a formal
proof of variance reduction under causal credit selection, a deterministic
safety bound for agentic rollout collection, and a four-factor experimental
protocol designed to produce falsifiable evidence for the central claims of
this work. We have further shown, via direct wall-clock instrumentation
across two independent 200-iteration training configurations, that the
infrastructure robustness guarantee of Theorem~\ref{thm:safety} holds in
practice independent of hardware configuration, that a naive implementation
of the causal event detector is responsible for a 1596\% rollout-collection
overhead concentrated almost entirely (98.0\%) in state-snapshot diffing, and
that a zero-extra-call HTTP-method-based detector is extensionally equivalent
to the naive detector under the REST semantics of the target environment.
Preliminary results confirm the robustness claim and the overhead-attribution
claim; the credit assignment claim awaits the controlled ablation and
horizon-scaling experiments described in Section~6, together with a
re-measurement of rollout-collection cost under the corrected detector
specified in Section~\ref{sec:http-detector}.

% =====================================================================
\begin{thebibliography}{9}

\bibitem{rudder2019}
J.~A. Arjona-Medina, M.~Gillhofer, M.~Widrich, T.~Unterthiner, J.~Brandstetter,
and S.~Hochreiter.
\textit{RUDDER: Return Decomposition for Delayed Rewards}.
NeurIPS, 2019.

\bibitem{schulman2016gae}
J.~Schulman, P.~Moritz, S.~Levine, M.~Jordan, and P.~Abbeel.
\textit{High-Dimensional Continuous Control Using Generalized Advantage Estimation}.
ICLR, 2016.

\bibitem{ppo2017}
J.~Schulman, F.~Wolski, P.~Dhariwal, A.~Radford, and O.~Klimov.
\textit{Proximal Policy Optimization Algorithms}.
arXiv:1707.06347, 2017.

\bibitem{loop2025}
Apple ML Research.
\textit{LOOP: Leave-One-Out Policy Optimization for Long-Horizon Agentic Tasks}.
2025.

\bibitem{grpo2024}
DeepSeek-AI.
\textit{Group Relative Policy Optimization}.
2024.

\bibitem{her2017}
M.~Andrychowicz, F.~Wolski, A.~Ray, et al.
\textit{Hindsight Experience Replay}.
NeurIPS, 2017.

\bibitem{nec2017}
A.~Pritzel, B.~Uria, S.~Srinivasan, et al.
\textit{Neural Episodic Control}.
ICML, 2017.

\bibitem{appworld2024}
H.~Trivedi, T.~Khot, M.~Hartmann, R.~Manku, V.~Dong, E.~Li, S.~Gupta,
A.~Sabharwal, and N.~Balasubramanian.
\textit{AppWorld: A Controllable World of Apps and People for Benchmarking
Interactive Coding Agents}.
ACL, 2024.

\bibitem{qwen2024}
Qwen Team, Alibaba Group.
\textit{Qwen2.5 Technical Report}.
arXiv:2412.15115, 2024.

\bibitem{lora2021}
E.~J. Hu, Y.~Shen, P.~Wallis, Z.~Allen-Zhu, Y.~Li, S.~Wang, L.~Wang, and
W.~Chen.
\textit{LoRA: Low-Rank Adaptation of Large Language Models}.
arXiv:2106.09685, 2021.

\bibitem{adamw2019}
I.~Loshchilov and F.~Hutter.
\textit{Decoupled Weight Decay Regularization}.
ICLR, 2019.

\bibitem{vllm2023}
W.~Kwon, Z.~Li, S.~Zhuang, Y.~Sheng, L.~Zheng, C.~H. Yu, J.~Gonzalez, H.~Zhang,
and I.~Stoica.
\textit{Efficient Memory Management for Large Language Model Serving with
PagedAttention}.
SOSP, 2023.

\bibitem{fsdp2023}
Y.~Zhao, A.~Gu, R.~Varma, L.~Luo, C.~C. Huang, M.~Xu, L.~Wright, H.~Shojanazeri,
M.~Ott, S.~Shleifer, A.~Desmaison, C.~Balioglu, P.~Damania, B.~Nguyen,
G.~Chauhan, Y.~Hao, A.~Mathews, and S.~Li.
\textit{PyTorch FSDP: Experiences on Scaling Fully Sharded Data Parallel}.
VLDB, 2023.

\end{thebibliography}

\end{document}
